\documentclass[11pt]{article}
\usepackage[margin=0.75in]{geometry}
\usepackage{amsmath,amssymb,siunitx}
\usepackage{parskip}
\usepackage{enumitem}

\title{PS160 --- Midterm 2 Equation Sheet}
\author{Thermodynamics \\ (Modules 17, 18, 19, 20)}
\date{}

\begin{document}
\maketitle

\section*{Temperature and Thermal Expansion (Module 17)}

\paragraph{Temperature scale conversions.}
\[
T_{K} = T_{C} + 273.15
\qquad
T_{F} = \tfrac{9}{5}T_{C} + 32
\qquad
T_{R} = T_{F} + 459.67
\]

\paragraph{Linear thermal expansion.}
\[
\Delta L = \alpha L_{0}\,\Delta T
\]

\paragraph{Volume thermal expansion.}
\[
\Delta V = \beta V_{0}\,\Delta T
\qquad
\beta = 3\alpha
\]

\paragraph{Thermal stress (rod prevented from expanding).}
\[
\frac{F}{A} = -Y \alpha\,\Delta T
\]

\paragraph{Quantity of heat / specific heat.}
\[
Q = m c\,\Delta T = n C\,\Delta T
\]

\paragraph{Latent heat (phase change).}
\[
Q = \pm m L
\qquad (L_{f}\text{: fusion},\; L_{v}\text{: vaporization})
\]

\paragraph{Heat transfer by conduction.}
\[
H = \frac{dQ}{dt} = k A\,\frac{\Delta T}{L}
\]

\paragraph{Heat transfer by radiation (Stefan--Boltzmann for a blackbody).}
\[
H = A\sigma T^{4}
\qquad\text{(emissivity $e$: } H = e A\sigma T^{4}\text{)}
\]
\[
\sigma = \SI{5.67e-8}{W/(m^{2} K^{4})}
\]

\paragraph{Useful constants.} $c_{\text{water}} = \SI{4190}{J/(kg\cdot K)}$.

\section*{Thermal Properties of Matter / Kinetic Theory (Module 18)}

\paragraph{Ideal gas law.}
\[
pV = N k T
\qquad
pV = n R T
\qquad
k N_{A} = R
\]

\paragraph{Boltzmann constant and Avogadro's number.}
\[
k = \SI{1.380649e-23}{J/K}
\qquad
N_{A} = \SI{6.02214076e23}{mol^{-1}}
\]

\paragraph{Equipartition theorem (each quadratic degree of freedom).}
\[
\langle \tfrac{1}{2} m v^{2}\rangle_{\text{per d.o.f.}} = \tfrac{1}{2} k T
\qquad
\tfrac{1}{2} m\,\langle v^{2}\rangle = \tfrac{3}{2} k T
\]

\paragraph{rms speed.}
\[
v_{\text{rms}} = \sqrt{\langle v^{2}\rangle} = \sqrt{\frac{3 k T}{m}} = \sqrt{\frac{3 R T}{M}}
\]

\paragraph{Maxwell--Boltzmann speed distribution.}
\[
f(v) \propto v^{2}\, e^{-m v^{2}/(2 k T)}
\]

\paragraph{Molar heat capacities for an ideal gas.}
\[
C_{V} = \frac{f}{2} R
\qquad
C_{P} = C_{V} + R
\qquad
\gamma = \frac{C_{P}}{C_{V}}
\]
Dulong--Petit for solids: $C_{V} = 3R$.

\section*{First Law of Thermodynamics (Module 19)}

\paragraph{Work done by a gas.}
\[
W = \int p\,dV
\]

\paragraph{First law.}
\[
dU = \bar{d}Q - \bar{d}W
\qquad
\Delta U = Q - W
\]

\paragraph{Internal energy of ideal gas (depends only on $T$).}
\[
U = n C_{V} T
\qquad
\Delta U = n C_{V}\,\Delta T
\]

\paragraph{Basic processes (ideal gas).}

Isobaric ($p$ const): $W = p\,\Delta V$, $Q = n C_{P}\,\Delta T$.

Isochoric ($V$ const): $W = 0$, $Q = n C_{V}\,\Delta T$, $\Delta U = Q$.

Isothermal ($T$ const): $\Delta U = 0$, $Q = W = n R T\ln(V_{2}/V_{1})$.

Adiabatic ($Q = 0$): $\Delta U = -W$ and
\[
p V^{\gamma} = \text{const},\quad
T V^{\gamma - 1} = \text{const},\quad
W = \frac{p_{1} V_{1} - p_{2} V_{2}}{\gamma - 1}
\]

\section*{Second Law / Entropy / Cycles (Module 20)}

\paragraph{Zeroth law.} $T_{A} = T_{B} = T_{C}$ (transitivity of thermal equilibrium).

\paragraph{Second law (Clausius / Kelvin--Planck).} $\Delta S_{\text{universe}} \ge 0$.

\paragraph{Entropy (definitions).}
\[
\Delta S = \int \frac{\bar{d}Q_{\text{rev}}}{T}
\qquad
S = k\ln w
\]

\paragraph{First law for a complete cycle.} $\Delta U = 0$, hence
\[
W_{\text{net}} = Q_{\text{net}} = Q_{H} + Q_{C}
\]
(with sign convention: $Q_{H} > 0$, $Q_{C} < 0$ for a heat engine).

\paragraph{Heat-engine efficiency.}
\[
e = \frac{W}{Q_{H}} = 1 + \frac{Q_{C}}{Q_{H}} = 1 - \frac{|Q_{C}|}{Q_{H}}
\]

\paragraph{Carnot efficiency (reversible maximum).}
\[
e_{\text{Carnot}} = 1 - \frac{T_{C}}{T_{H}}
\]

\paragraph{Refrigerator coefficient of performance.}
\[
K_{\text{ref}} = \frac{Q_{C}}{-W} = \frac{|Q_{C}|}{|W|}
\qquad
K_{\text{heat pump}} = \frac{-Q_{H}}{-W} = \frac{|Q_{H}|}{|W|}
\]

\paragraph{Carnot cycle entropy change.}
\[
\Delta S_{\text{Carnot}} = \frac{Q_{H}}{T_{H}} + \frac{Q_{C}}{T_{C}} = 0
\]

\paragraph{Entropy of simple processes.}
Isothermal ideal gas: $\Delta S = n R\ln(V_{2}/V_{1})$.\\
Heating a substance: $\Delta S = m c\ln(T_{2}/T_{1})$.\\
Phase change: $\Delta S = \pm m L / T$.

\end{document}
